% Documento LaTeX para Mecânica das Rochas - Versão Simplificada
\documentclass[12pt,a4paper,twoside]{book}
\usepackage[portuguese]{babel}
\usepackage{geometry}
\geometry{margin=2.5cm, headheight=15pt}

\usepackage{graphicx}
\usepackage{pdfpages}
\usepackage{fancyhdr}
\usepackage{hyperref}
\usepackage{xcolor}
\usepackage{amsmath}
\usepackage{amssymb}

% Configuracao de cores
\definecolor{darkblue}{rgb}{0.1,0.1,0.5}
\hypersetup{colorlinks=true, linkcolor=darkblue, urlcolor=darkblue}

% Cabecalho e rodape
\pagestyle{fancy}
\fancyhf{}
\fancyhead[LE,RO]{\rightmark}
\fancyhead[LO,RE]{\leftmark}
\fancyfoot[C]{\thepage}
\renewcommand{\headrulewidth}{0.5pt}
\setlength{\headheight}{15.04742pt}

% Titulo
\title{\huge \textbf{Mecanica das Rochas} \\
        \large Curso de Licenciatura em Engenharia de Petroles \\
        \normalsize 3.º Ano}
\author{Autor: Felisberto Queta}
\date{\today}

\begin{document}

% Capa
\maketitle

% Indice
\tableofcontents
\newpage

% ============================================================================
% CAPITULO 1 - INTRODUCAO A MECANICA DAS ROCHAS
% ============================================================================
\chapter{Introducao a Mecanica das Rochas}

\section{Definicao e Conceitos Fundamentais}

\textbf{Mecanica das Rochas} e uma disciplina da engenharia civil que se ocupa do estudo teorico 
e pratico das propriedades das rochas e das suas respostas quando submetidas a uma perturbacao 
aplicada. Estas respostas sao chamadas deformacoes que se podem manifestar em forma de:

\begin{itemize}
    \item \textbf{Juntas ou diaclases}: Fracturacao sem deslocamento percetivel dos blocos.
    \item \textbf{Falhas}: Torcao com deslocamento relativo de blocos um em relacao ao plano de fractura.
    \item \textbf{Dobras}: Flexao, encurtamento acentuado e deformacao interna da rocha.
\end{itemize}

\section{O Papel do Engenheiro de Petroleo}

Durante suas atividades profissionais, o engenheiro de petroleo lida com:
\begin{itemize}
    \item Exploracao e producao de hidrocarbonetos
    \item Gestao da comercializacao e transporte de hidrocarbonetos
    \item Prospeccao e pesquisa de bacias sedimentares com potencial economico
    \item Acompanhamento da producao de derivados do petroleo
\end{itemize}

Para isso, necessita de conhecimentos em:
\begin{itemize}
    \item Geologia
    \item Geofisica
    \item Petrofisica
    \item \textbf{Mecanica das Rochas}
\end{itemize}

\subsection{O Conceito de Geomecanica}

\textbf{Geomecanica} e a disciplina que procura integrar Geologia, Geofisica, Petrofisica e 
Mecanica das Rochas, munindo o engenheiro de petroleo de habilidades e competencias para melhor 
compreender suas atividades e as de seus colegas em outras areas.

\subsection{Consequencias da Falta de Dominio em Geomecanica}

Quando o engenheiro nao domina adequadamente os conceitos de geomecanica, podem ocorrer:
\begin{itemize}
    \item Desmoronamento das paredes do poco
    \item Explosao de plataforma (blowout) devido ao vazamento de fluidos
    \item Descoberta de pocos secos
\end{itemize}

\section{Objetivos Gerais da Disciplina}

\begin{enumerate}
    \item Compreender os principais conceitos de Mecanica das Rochas e sua importancia nos varios 
    setores, principalmente na industria petrolífera.
    \item Identificar os diferentes tipos de rochas e diferencia-las de acordo com suas caracteristicas.
    \item Determinar as propriedades fisicas e mecanicas de uma rocha.
    \item Compreender e aplicar as nocoes sobre estado de tensao, resistencia e deformabilidade 
    nos macos rochosos.
    \item Perceber os tipos de descontinuidades e sua importancia no comportamento de macos rochosos.
    \item Perceber a importancia do estudo do comportamento das rochas no entendimento das estruturas 
    da Terra, na economia dos paises e na preservacao do ambiente.
    \item Realizar classificacao geomecanica e descricao dos macos rochosos.
\end{enumerate}

\section{Aplicacoes na Industria Petrolifera}

\subsection{1. Analise de Bacias}

Compreende o estudo da:
\begin{itemize}
    \item Origem das bacias
    \item Evolucao das bacias
    \item Inversao de bacias
\end{itemize}

\subsection{2. Perfuracao de Pocos}

E necessario:
\begin{itemize}
    \item Analise de tensoes in situ para prever colapsos e perdas de fluido
    \item Determinacao da pressao ideal da lama de perfuracao
    \item Modelagem do comportamento geomecanico do poco
\end{itemize}

\section{Aplicacoes na Engenharia Civil}

\subsection{Tuneis e Galerias Subterraneas}

O estudo da estabilidade de macos rochosos e essencial para a construcao de:
\begin{itemize}
    \item Tuneis rodoviarios
    \item Tuneis ferroviarios
    \item Tuneis metroviarios
\end{itemize}

\subsection{Fundacoes de Grandes Estruturas}

Projetos de barragens, pontes e edificios exigem analise detalhada das condicoes do maco rochoso 
para garantir seguranca.

\subsection{Taludes e Encostas}

A estabilizacao de encostas e cortes rochosos e essencial para prevenir deslizamentos em:
\begin{itemize}
    \item Rodovias
    \item Ferrovias
    \item Terrenos em geral
\end{itemize}

\section{Aplicacoes na Engenharia de Minas}

A Mecanica das Rochas permite compreender o comportamento das rochas em diferentes condicoes de 
carregamento, contribuindo para a seguranca e eficiencia das operacoes mineiras.

\subsection{Areas Principais de Aplicacao}

\begin{enumerate}
    \item Estabilidade de Escavacoes
    \item Dimensionamento de Suporte
    \item Fragmentacao e Perfuracao de Rochas
    \item Controlo de Pressoes e Deformacoes
    \item Drenagem e Controlo de Agua
    \item Planeamento e Recuperacao
\end{enumerate}

% ============================================================================
% CAPITULO 2 - ROCHAS, CONCEITOS E CARACTERISTICAS
% ============================================================================
\chapter{Rochas, Conceitos e Caracteristicas}

\section{Definicao de Rochas}

\textbf{Rochas} sao materiais solidos, consolidados e constituidos naturalmente por um ou mais 
minerais, com caracteristicas fisicas e mecanicas especificas para cada tipo.

\section{Formacao das Rochas}

As rochas formam-se atraves de diferentes processos da geodinamica interna e externa:
\begin{itemize}
    \item \textbf{Magmatismo}: Formacao de rochas igneas ou magmaticas
    \item \textbf{Sedimentacao e Diagenese}: Formacao de rochas sedimentares
    \item \textbf{Metamorfismo}: Formacao de rochas metamorficas
\end{itemize}

\section{Tipos de Rochas}

As rochas podem classificar-se em tres grupos principais baseado na sua composicao, textura e 
caracteristicas geneticas.

\subsection{Rochas Igneas (Magmaticas)}

Formam-se a partir da cristalizacao do magma e constituem cerca de 20\% de todas as rochas.

\subsubsection{Rochas Plutonicas (Intrusivas)}

Quando o magma arrefece-se lentamente em profundidade:
\begin{itemize}
    \item \textbf{Textura}: Faneritica (graos grandes e bem formados)
    \item \textbf{Exemplos}: Granitos, Gabros, Sienitos, Dioritos
\end{itemize}

\subsubsection{Rochas Vulcanicas (Extrusivas)}

Quando o magma arrefece-se rapidamente a superfície:
\begin{itemize}
    \item \textbf{Textura Vitrea}: Material vitreo sem cristais (ex: Obsidiana)
    \item \textbf{Textura Afanitica}: Cristais muito finos (ex: Riólitos, Basaltos)
\end{itemize}

\subsection{Rochas Metamorficas}

Resultam da transformacao mecanica, termica e quimica de rochas pre-existentes. Constituem cerca 
de 14\% de todas as rochas.

\subsubsection{Tipos de Metamorfismo}

\begin{itemize}
    \item \textbf{Metamorfismo Regional}: Abrange extensas areas, particularmente em zonas de 
    formacao de cadeias de montanha.
    
    \item \textbf{Metamorfismo de Contacto}: Limitado ao contacto entre a litologia presente 
    e a intrusao de magma.
\end{itemize}

\subsection{Rochas Sedimentares}

Constituem cerca de 66\% de todas as rochas da crusta terrestre. Resultam do transporte, deposicao 
e consolidacao de sedimentos.

\subsubsection{Classificacao das Rochas Sedimentares}

\begin{enumerate}
    \item \textbf{Rochas Sedimentares Clasticas}: Resultam do transporte, deposicao e consolidacao 
    de sedimentos provenientes da erosao de rochas pre-existentes.
    
    \item \textbf{Rochas Sedimentares Quimiogenicas}: Formadas pela precipitacao quimica de minerais 
    da agua.
    
    \item \textbf{Rochas Sedimentares Biogenicas}: Formam-se a partir da acumulacao de materia organica.
\end{enumerate}

\subsubsection{Importancia na Industria de Petroleo}

\begin{itemize}
    \item \textbf{Arenito}: Rocha-reservatorio preferida (alta porosidade)
    \item \textbf{Folhelho}: Rocha geradora de petroleo
    \item \textbf{Evaporito}: Rocha selante
    \item \textbf{Calcario}: Rocha geradora e reservatorio
\end{itemize}

\section{Caracteristicas Principais das Rochas por Tipo}

\begin{table}[h]
\centering
\begin{tabular}{|l|c|c|c|}
\hline
\textbf{Caracteristica} & \textbf{Magmaticas} & \textbf{Metamorficas} & \textbf{Sedimentares} \\
\hline
\textbf{Aspecto} & Macico e compacto & Macico ou foliado & Em camadas, friáveis \\
\hline
\textbf{Textura} & Graos interconectados & Graos interconectados & Graos nao imbricados \\
& sem orientacao & com orientacao & com poros \\
\hline
\textbf{Forma} & Formas irregulares & Formas irregulares & Arredondadas/ovuladas \\
\hline
\end{tabular}
\caption{Caracteristicas principais das rochas}
\label{tab:rochas}
\end{table}

\section{Solos}

\subsection{Definicao}

\textbf{Solos} sao agregados de graos minerais nao consolidados encontrados na superfície terrestre 
que nao oferecem resistencia a escavacao mecanica.

\subsection{Classificacao de Solos}

\begin{itemize}
    \item \textbf{Solos Residuais}: Resultantes da meteoricao da rocha sem transporte
    \item \textbf{Solos Sedimentares}: Resultantes de meteoricao, transporte e deposicao
\end{itemize}

\subsection{Diferencas entre Rochas e Solos}

\begin{table}[h]
\centering
\begin{tabular}{|l|l|}
\hline
\textbf{Rochas} & \textbf{Solos} \\
\hline
Modulos de deformacao altos & Modulos de deformacao baixos \\
\hline
Elevada coesao interna & Coesao baixa ou nula \\
\hline
Elevados valores de resistencia a tracao & Baixa resistencia a tracao \\
\hline
Baixa permeabilidade & Alta permeabilidade \\
\hline
\end{tabular}
\caption{Diferencas entre rochas e solos}
\label{tab:rochas_solos}
\end{table}

\section{Macos Rochosos}

\subsection{Definicao}

Um \textbf{maco rochoso} e uma porcao definida de uma ou mais formacoes geologicas, caracterizada 
por suas rochas e descontinuidades.

\subsection{Caracteristicas Principais}

\begin{itemize}
    \item \textbf{Descontinuidade}: Qualquer plano que separa os blocos da matriz rochosa.
    \item \textbf{Anisotropia}: Comportamento diferente conforme a direcao considerada.
    \item \textbf{Heterogeneidade}: Materiais com caracteristicas distintas.
    \item \textbf{Matriz Rochosa}: Material isento de descontinuidades.
\end{itemize}

\subsection{Tipos de Macos}

\begin{itemize}
    \item \textbf{Maco Intacto}: Continuo, isotropico, homogeneo
    \item \textbf{Maco Descontinuo}: Fracturado, anisotropico, heterogeneo
\end{itemize}

\subsection{Fatores Geologicos Determinantes}

\begin{enumerate}
    \item Litologia e Propriedades da Matriz Rochosa
    \item Grau de Alteracao ou Meteoricao
    \item Estructura Geologica e Descontinuidades
    \item Estado de Esforcos
    \item Condicoes Hidrogeologicas
\end{enumerate}

\newpage

% ============================================================================
% ANEXOS - INCLUIR PDFs
% ============================================================================
\appendix
\chapter{Documentos Originais}

Os documentos originais em formato PDF estao incluidos abaixo como referencia visual.

\section{Apresentacao e Introducao a Mecanica das Rochas 2025/2026}


\includepdf[pages=-,landscape=false,scale=0.85]{referências para o estudo/1. Apresentação e Introdução a Mecânica das Rochas 2025_2026 (Parte 1).pdf}

\section{Introducao a Mecanica das Rochas 2025 (Completo)}

\includepdf[pages=-,landscape=false,scale=0.85]{referências para o estudo/1. Introdução à Mecânica das Rochas 2025_Completo.pdf}

\section{Rochas, Conceitos e Caracteristicas 2025}

\includepdf[pages=-,landscape=false,scale=0.85]{referências para o estudo/2. Rochas, conceitos e características_MR 1_2025.pdf}

\section{Apresentacao e Introducao 2024}


\includepdf[pages=-,landscape=false,scale=0.85]{referências para o estudo/Apresentação e Introdução da Mecânica das Rochas 2024.pdf}

\section{Apresentacao e Introducao 2024 (Duplicado)}


\includepdf[pages=-,landscape=false,scale=0.85]{referências para o estudo/Apresentação e Introdução da Mecânica das Rochas 2024 (1).pdf}

\end{document}
